\chapter{Discussione}
\label{chap:discussione}

Le teste ibride ottengono buone prestazioni sulla rappresentazione
estratta da ResNet18, ma non superano il controllo classico B2 nei
confronti statistici del \cref{chap:risultati}. Per interpretare questo
risultato occorre considerare separatamente la difficoltà del task,
la disponibilità dei dati e il costo. La discussione segue le domande
di ricerca del \cref{sec:intro-rq}; le possibili spiegazioni sono
distinte dai meccanismi effettivamente verificati.

\section{Confronto principale e ruolo della rappresentazione}
\label{sec:disc-rq1}

Per RQ1, i dati non forniscono evidenza di un miglioramento di Q
rispetto a B2. Sul Full binario
le differenze medie sono piccole, pari a \(+0{,}02\) e \(+0{,}19\)
punti percentuali per quattro e otto qubit, senza significatività
dopo Holm. Sul Full a tre classi si osservano invece differenze di
\(-0{,}73\) e \(-1{,}67\) pp, significative secondo la procedura
adottata. L'assenza di significatività nel primo caso non dimostra
equivalenza, che richiederebbe un margine definito e un test dedicato.

Le buone prestazioni di B0 e B1 aiutano a interpretare questo quadro.
Una parte rilevante dell'informazione discriminativa è già contenuta
nelle feature classiche. La componente quantistica non deve apprendere
una rappresentazione direttamente dai pixel, ma rielaborare una
rappresentazione preesistente dopo una forte compressione. La domanda
non riguarda pertanto la superiorità di un intero sistema di visione
quantistico, bensì il valore aggiunto di una particolare testa.

B2 conserva prestazioni competitive con un bottleneck di sole quattro
o otto componenti. Ciò indica che, nelle condizioni osservate, la
compressione non impedisce di ottenere buoni risultati. Non dimostra
che le altre componenti delle feature siano sempre inutili, né che
qualsiasi proiezione verso poche dimensioni sarebbe sufficiente:
la proiezione è addestrata insieme al classificatore.

Nel binario il livello prossimo a \(0{,}96\) di Macro-F1 è compatibile
con una situazione in cui il margine di miglioramento misurabile è
contenuto. L'espressione effetto soffitto va tuttavia intesa come
interpretazione: non sono noti l'errore irriducibile del problema o
il massimo conseguibile con le etichette disponibili. La convergenza
delle medie non basta a identificarli.

\section{Complessità del task e classe merger}
\label{sec:disc-rq2}

La risposta a RQ2 è che il divario relativo diventa più sfavorevole
a Q nel passaggio alle tre classi. La differenza-in-differenze sul
Full è di circa \(-0{,}75\) pp a quattro qubit e \(-1{,}86\) pp
a otto. Questo contrasto è più informativo del solo calo assoluto
di Macro-F1, che interessa tutti i modelli.

Non si tratta però di un effetto causale isolato dell'aggiunta dei
merger. Cambiano il numero di classi, la composizione degli esempi
e, sul Full, la numerosità del training. Inoltre, la Macro-F1
media due oppure tre F1 per classe. La conclusione sostenibile è
quindi una variazione osservata del confronto relativo tra task,
non la quantificazione di una penalità fisica attribuibile a una
sola categoria.

I sistemi etichettati come merger costituiscono un banco di prova utile:
formano la classe con F1 media più bassa sul Full per B2, Q e QF. A otto qubit,
il loro contributo aritmetico rappresenta poco più della metà dello
scarto complessivo Q--B2 nella Macro-F1. Le matrici di confusione
mostrano anche più merger assegnate alle spirali da Q-8. A quattro
qubit, invece, il quadro include una perdita di precision merger
con recall lievemente superiore al controllo. Non esiste quindi
un'unica modalità di errore comune alle due configurazioni.

La maggiore eterogeneità visuale dei sistemi perturbati è una
spiegazione plausibile della difficoltà condivisa dai modelli.
I risultati non permettono però di distinguere quanto dipenda dalla
rappresentazione congelata, dalla testa, dall'ambiguità delle etichette
o dalle condizioni osservative. Attribuire gli errori direttamente
all'incapacità del circuito di riconoscere specifiche strutture
astronomiche richiederebbe esperimenti ulteriori.

\section{Disponibilità dei dati e regimi low-resource}
\label{sec:disc-rq3}

RQ3 verifica se il circuito sia particolarmente utile quando il
training contiene pochi esempi. Le evidenze non sostengono questa
ipotesi nel protocollo considerato. Nelle tre classi Q ha una media
inferiore a B2 in tutti i regimi; a otto qubit le differenze dei
tre subset rimangono significative dopo la correzione secondaria.
Il divario medio passa da \(-3{,}30\) pp con 100 immagini a
\(-1{,}67\) pp sul Full: cresce la prestazione di entrambi e lo
scarto si riduce, senza annullarsi.

Questo andamento mostra perché sia utile non limitarsi al training
completo. Le prestazioni vicine sul task binario non descrivono
l'intero comportamento delle architetture; task più articolati e
subset ridotti possono rendere visibili differenze attenuate altrove.
La F1 merger con 100 immagini e otto qubit, inferiore in Q di circa
\(9{,}82\) pp rispetto a B2, rappresenta uno di questi casi,
da interpretare come analisi esplorativa di una singola classe.

La selezione dei subset dal medesimo training mantiene controllata
la provenienza dei dati, senza dimostrare l'annidamento reciproco
dei campioni ridotti. Le variazioni possono quindi riflettere anche
la composizione dei campioni. Inoltre, con batch e limite di epoche
fissati, un dataset più piccolo produce meno aggiornamenti per epoca;
l'early stopping modifica inoltre la durata effettiva. Le curve
descrivono quindi il comportamento del protocollo a differenti
numerosità, non un esperimento a identico numero di aggiornamenti.

\section{Addestramento circuitale e numero di qubit}
\label{sec:disc-rq4}

Il confronto Q--QF risponde alla prima parte di RQ4. Sul Full i
vantaggi medi di Q sono compresi tra circa \(0{,}10\) e \(0{,}26\)
pp e tutti gli intervalli includono lo zero. Non emerge dunque un
beneficio netto e uniforme dell'ottimizzazione circuitale, pur in
presenza di alcuni effetti esplorativi più favorevoli nei subset.

La vicinanza tra Q e QF è compatibile con un ruolo importante degli
strati classici che precedono e seguono il circuito. QF può adattare
gli angoli di ingresso attraverso la proiezione e modificare la
decisione finale attraverso il classificatore: non è un modello
privo di apprendimento. Non si può quindi dedurre che il circuito
sia ignorato, che i suoi gradienti siano nulli o che la rete sia
equivalente a un classificatore lineare.

La seconda parte di RQ4 riguarda la dimensionalità. Otto qubit non
offrono un miglioramento sistematico: nel task a tre classi Q-8
ha una media inferiore a Q-4 in tutti i regimi, mentre nel binario
con 100 esempi l'ordinamento è opposto. Più qubit non costituiscono
quindi, in queste configurazioni, una garanzia di maggiore utilità
predittiva.

Il passaggio da quattro a otto modifica però anche proiezione,
classificatore, numero di rotazioni e collegamenti circuitali.
Il confronto non identifica l'effetto separato di ciascun elemento.
Una maggiore difficoltà di ottimizzazione o una diversa adeguatezza
della trasformazione ai dati sono ipotesi possibili, non meccanismi
dimostrati dal solo ordinamento delle medie.

In particolare, non viene diagnosticato un barren plateau. Il fenomeno
descritto da McClean et al.~\cite{mcclean2018barren} richiede condizioni
e misure specifiche. Come chiarito nel metodo, il campo diagnostico
denominato \texttt{quantum\_gradient\_norm} include anche parametri
classici e non permette di isolare i gradienti degli angoli variazionali.
Né la loro scomparsa né la loro adeguata ampiezza possono essere
concluse da quella diagnostica aggregata.

\section{Qualità probabilistica e costo computazionale}
\label{sec:disc-rq5}

Le metriche probabilistiche aggiungono informazioni che la Macro-F1
non contiene. Un modello può avere decisioni meno accurate e una
confidenza mediamente più coerente con la propria frequenza di successo.
Questo aiuta a leggere il caso di Q-8 sulle tre classi: l'ECE medio
è inferiore a quello di B2-8, ma Brier score e log-loss sono peggiori.
Non è una contraddizione e non costituisce un ribaltamento del
confronto predittivo. Le metriche valutano aspetti diversi e l'ECE
dipende anche dal binning \cite{guo2017calibration,gneiting2007scoring}.

Per RQ5 il costo osservato è nettamente sfavorevole ai modelli
quantistici simulati. Sul Full, le mediane dei rapporti appaiati
di training Q/B2 sono circa dieci a quattro qubit e diciannove a
otto. Nelle tre classi B2 combina inoltre una Macro-F1 media
superiore con tempi mediani inferiori. Nel binario i piccoli
vantaggi medi di alcune configurazioni Q non sono statisticamente
supportati e richiedono comunque un costo maggiore.

QF riduce il costo di training rispetto a Q sul Full, ma resta
più costoso di B2. Congelare il circuito non evita la sua esecuzione
né la differenziazione rispetto agli ingressi necessaria per
addestrare la proiezione. La distinzione tra pesi congelati e
assenza di calcolo del gradiente è quindi rilevante anche dal
punto di vista computazionale.

I tempi riguardano le teste simulate ed escludono l'estrazione delle
feature. Il diverso ambiente di esecuzione tra regimi
(\cref{sec:ambiente-hardware}) e il numero di epoche determinato
dall'early stopping impediscono di attribuire tutte le variazioni
temporali alla numerosità dei dati o al costo delle porte.
Queste misure non consentono inoltre di stimare tempi o consumi di una QPU.

\section{Relazione con la letteratura}

Il lavoro conferma la praticabilità dello schema di transfer learning
ibrido introdotto da Mari et al.~\cite{mari2020transfer}: la
rappresentazione classica consente di utilizzare piccoli circuiti
su immagini di dimensione elevata. Questa fattibilità è distinta
dalla superiorità predittiva. Anche l'interesse per piccoli dataset
documentato nel telerilevamento \cite{otgonbaatar2023transfer} motiva
una verifica, ma non implica che lo stesso comportamento debba
emergere sulle galassie.

I risultati sono coerenti con l'esigenza metodologica evidenziata
da Bowles et al.~\cite{bowles2024benchmarking}: la scelta di
baseline e condizioni sperimentali influenza ciò che si può
attribuire alla componente quantistica. Qui B2 è particolarmente
informativo perché conserva il bottleneck e gli strati classici
di interfaccia senza inserire il circuito.

Questa coerenza non rende equivalenti gli esperimenti. Dataset,
encoding, ansatz, ottimizzatori e criteri di selezione differiscono
tra gli studi. Il contributo della tesi è un'evidenza documentata
su un preciso problema applicativo, non una confutazione generale
delle possibilità del quantum machine learning.

\section{Limiti e minacce alla validità}
\label{sec:disc-limiti}

\subsection{Campione, etichette e generalizzazione esterna}

Il dataset è un campione selezionato e bilanciato, non una rappresentazione
delle frequenze morfologiche di una survey completa. Le soglie sui voti
definiscono classi operative e possono privilegiare esempi con maggiore
accordo. Le metriche ottenute non determinano quindi direttamente
le prestazioni su popolazioni con prevalenze, qualità osservative o
livelli di ambiguità diversi.

Le partizioni separano gli identificativi degli oggetti, ma non sono
costruite raggruppando sistemi fisici o regioni di cielo. L'unicità
degli identificativi non garantisce l'indipendenza astronomica di
tutti gli esempi. Una validazione esterna e controlli sulle eventuali
correlazioni spaziali rafforzerebbero la valutazione della generalizzazione.

\subsection{Rappresentazione, iperparametri e ambito architetturale}

Il backbone congelato rende il confronto più controllato, ma limita tutte le teste
all'informazione disponibile nelle feature estratte. Non viene studiato
se il fine-tuning sul dominio galattico cambi il confronto Q--B2.
I risultati riguardano angle encoding singolo, profondità due e
quattro oppure otto qubit; non includono altre famiglie di circuiti,
rumore fisico o misure a shot finiti.

Gli iperparametri fissati favoriscono la comparabilità ma non equivalgono
a una ricerca esaustiva del migliore assetto per ogni architettura.
Una configurazione condivisa può essere più adatta a un modello
che a un altro. B0 e B1 usano inoltre batch diversi dal confronto
principale: le loro prestazioni offrono contesto, non isolano
esclusivamente l'effetto della forma della testa.

\subsection{Ripetizioni e inferenza statistica}

Le dieci ripetizioni condividono il dataset e possono includere
gli stessi oggetti in test differenti. Non sono dieci campioni
astronomici indipendenti. Il bootstrap sui seed e il test sign-flip
non modellano esplicitamente questa dipendenza, mentre la correzione
di Holm tratta la molteplicità e non corregge le assunzioni dei test
di partenza. Le cautele discusse da Nadeau e Bengio
\cite{nadeau2003inference} restano quindi pertinenti.

L'appaiamento riduce la confusione dovuta a confronti su split
diversi, ma non elimina questi limiti. Intervalli, p-value e
dimensioni dell'effetto devono essere letti insieme, distinguendo
evidenze primarie e osservazioni esplorative. Il numero complessivo
di 640 run non va scambiato per 640 repliche indipendenti di
ciascun confronto: ogni contrasto utilizza dieci coppie.

\subsection{Audit, casi particolari e interpretabilità}

La conservazione di predizioni, checkpoint e cronologie consente
di controllare la coerenza degli aggregati e di esaminare casi
particolari. Un seed favorevole o un'immagine classificata
correttamente soltanto da Q può suggerire un'analisi, ma non
sostituisce il confronto aggregato e non giustifica l'esclusione
di run meno favorevoli.

Analogamente, una mappa di attribuzione visuale evidenzia contributi
locali alla risposta del modello secondo il metodo utilizzato,
non una spiegazione causale completa della decisione. Le analisi
qualitative devono essere accompagnate da criteri espliciti di
selezione degli esempi e non usate per trasformare un caso
illustrativo in una proprietà generale dell'architettura.

\section{Sintesi interpretativa}

Il quadro osservato distingue la capacità di classificare dall'utilità
di introdurre il circuito. Le teste ibride raggiungono prestazioni
elevate sul binario, ma il controllo classico ottiene risultati
vicini con costi inferiori. Nelle tre classi e soprattutto in
alcuni regimi ridotti il divario diventa più sfavorevole a Q.
L'addestramento circuitale e l'aumento dei qubit non offrono
benefici uniformi.

Per le configurazioni studiate, il costo aggiuntivo non è accompagnato
da un vantaggio predittivo dimostrato. Una prosecuzione motivata del
lavoro è verificare quali trasformazioni siano adatte alle feature
disponibili e quanto l'ottimizzazione riesca a sfruttarle, prima di
aumentare ulteriormente la complessità del circuito. È una direzione
di ricerca suggerita dagli esperimenti, non una conclusione generale
sull'utilità del quantum machine learning.
